\documentclass[12pt]{article}
\usepackage{graphicx} % Required for inserting images
\usepackage{newtx}
\usepackage[margin=1in]{geometry}
\usepackage{indentfirst}
\usepackage{url}

\usepackage{titlesec}
\titleformat{\section}
  {\normalfont\fontsize{12}{12}\selectfont\bfseries}
  {\thesection}{1em}{}

\setcounter{secnumdepth}{0}

\usepackage{setspace}
\doublespacing % Applies double spacing globally

\date{}

\title{
  \fontsize{14}{16}\selectfont \underline{
    \textbf{Steering Optical Signals In The Visible Light Spectrum Using Beamforming}
    }
}

\begin{document}

\maketitle
\vspace{-7em}

\section{Background}
\vspace{-1em}
\par
Many modern systems in many different fields involve steering
electromagnetic wavefronts at various different wavelengths. At smaller
wavelengths, the techniques for steering these electromagnetic waves
fall under the category of optics, and the steering is typically carried
out with specialized tools called lenses, which are specialized
materials in specific shapes that affect the velocity of the wave moving
through in such a way as to change the direction of the wave
propagation. In many circumstances, lenses can present their own
engineering challenges, such as heating under high energy, and can take
a lot of space in space-constrained applications.
\par
In larger wavelength electromagnetic wave applications, such as radar, and
in acoustic applications, there is another common technique for steering
the propagation of waves, called beamforming. Beamforming is a signal
processing technique that steers the wavefront by slightly altering the
phase of a given signal broadcast from multiple antennas in a single
array, causing constructive interference in the desired direction, and
destructive interference in every other direction \cite{Lichtman_PySDR_2023}.
\par
While beamforming is an established signal processing technique, and used
widely in many fields, in visible light applications it is mainly used for Visible
Light Communication \cite{GUPTA2024111182}. This application is useful, but low power
and limited in scope. More general applications, like replacing optical systems, have
had not had very much exploration.
\par
This project aims to apply beamforming techniques to
visible light, in order to steer the visible light in a way similar to
optical lenses, to attempt to create an optical system that does not
require physical lenses.

\section{Question}
Can beamforming techniques be used to steer visible light signals,
and can this be used as a suitable alternative to physical optics?

\section{Hypothesis}
It is possible to steer visible light signals using beamforming techniques, 
and this can be used as a suitable alternative to physical optics in some applications.

\section{Methods}
The first part of this experiment will involve developing computer simulations of the 
wave propagation of visible light. This will lead into simulating a physical source
array, and these simulations will be used to develop a model for a physical array. Twenty
hours across the span of two weeks will be allotted for the development and validation of 
these simulation models. 
\par
Next, the physical array will be chosen and constructed. This will likely consist of some sort of
LED display, but the simulations will provide more exact array implementation details. The 
array construction will also include a light sensor array\cite{Analog_Devices}, to quantifiably validate the 
beamforming results. The construction and testing of this setup will take an additional 
twenty hours across two weeks. 
\par
The third step of this experiment will involve writing a program to turn the simulated
beamforming steering math into some instructions the array can actually execute. This 
will involve writing code to individually control subsections of the array, and to pulse
those sections with the frequency and phase shifts necessary for beamforming to take place.
Presumably this will be the most complicated part of the project, and fourty hours will be
allotted to this across the span of four weeks.
\par
The fourth step of this experiment involves quantifiably validating the experimental results,
and will involve writing a program to read from the aforementioned light sensor and log the 
recieved light data at different locations relative to the array. The point of this is to show
that beamforming is actually occurring, and that will be shown by an increase in brightness at 
the intended location. This will take an additional twenty hours over two weeks.
\par
Finally, the results will be analyzed, and a final report and presentation will be written to
summarize the findings of the experiment. Another twenty hours across two weeks will be allotted
for this.

\section{Significance}
The aim of this experiment is to determine if beamforming techniques could be used to supplant
traditional optical systems. If this is the casae, many optical applications can become
very simplified and much more efficient. This technique could be useful in heads up displays,
virtual reality, telescopes, and even Extreme UltraViolet (EUV) semiconductor manufacturing. The
benefits of this could be far reaching, reducing cost and leading to the development of new, simpler,
and more efficient technologies.


\bibliographystyle{unsrturl}
\bibliography{citations}

\end{document}
